% No preamble. Prefix fml-. Qualitative uniform finite leakage on cofinal periods.
\section{Uniform disappearance of finite classical marginal leakage}
\label{fml-section}

This section gives a qualitative uniform statement along the
standard nested compatible period family. It does not assert
exact finite marginal preservation or an exponential convergence
rate. The preceding classical depth may vary arbitrarily with
the period. The proof uses the already established exact
preservation of the true infinite Maxwell coefficient in each
joint classical system, rather than identifying different
limiting systems.

\subsection{Period family and statement}
\label{fml-setup-section}

Fix the primitive ratio \(\ell\), an admissible composite ratio
\(R=\ell^s\), auxiliary grouping \(J\), carrier exponent \(a_0\),
and one positive common analytic radius \(b\). All these data
are fixed. Let \(\mathcal P\) be a compatible nested family of
cubic output periods tending to infinity and satisfying
\begin{equation}
 \text{for each fixed }q\ge0,\quad
 R^q\mid N\quad\hbox{for all sufficiently late }N\in\mathcal P.
 \label{fml-divisibility}
\end{equation}
All grouping and source-geometry divisibilities are also
imposed. The standard nested source tori have this property;
for example \(N=2\ell^k\) does when the fixed grouping is
compatible. Arbitrary growing multiples of \(J\) with bounded
\(\ell\)-adic valuation are not covered by this argument.

Let \(\mathcal T_{M,R,N}\) be the actual pre-gauge classical
substitution from the input torus of current period \(RN\)
to the output torus of period \(N\). Use the actual finite
short-lift Maxwell ratios, each normalized by its own actual
finite homogeneous reference, and define
\begin{equation}
 \Lambda_{M,R,N}
   =\beta_{MR,N}\mathcal T_{M,R,N}-\beta_{M,RN}.
 \label{fml-leakage}
\end{equation}
This is a linear functional on the invariant finite activity
quotient at the common \(a_0,b\). On an extracted input it
is the finite classical marginal leakage.

\begin{theorem}[Qualitative uniform finite leakage]
\label{fml-theorem}
For all sufficiently large lower period thresholds \(K\), set
\begin{equation}
 \omega_{R,J,b}(K)
   =\sup_{\substack{M,\ N\in\mathcal P\\N\ge K}}
      \ \sup_{\substack{x\ {\rm invariant}\\\|x\|\le b^2}}
                     |\Lambda_{M,R,N}x|.
 \label{fml-modulus}
\end{equation}
Then
\begin{equation}
                         \omega_{R,J,b}(K)\longrightarrow0
                         \quad(K\longrightarrow\infty).
 \label{fml-vanishing}
\end{equation}
The supremum allows every permitted preceding homogeneous
depth \(M\), with no fixed-depth or bounded-depth restriction.
\end{theorem}

The function in \eqref{fml-modulus} is finite and decreasing
above the reference thresholds. Indeed the uniform finite
moment bound and fixed-\(R\) transport bound give
\begin{equation}
 \omega_{R,J,b}(K)\le C_\beta(B_R+1).
 \label{fml-crude-bound}
\end{equation}
The factor \(b^{-2}\) in the dual moment norm cancels the
input normalization \(b^2\). The theorem asserts a
qualitative improvement of this bound, not a rate.

\subsection{Compactness of varying finite input classes}
\label{fml-input-compactness}

Take any sequence \(M_n\), periods \(N_n\in\mathcal P\)
tending to infinity, and invariant finite input classes
\(x_n\) with \(\|x_n\|\le b^2\). Choose representatives
\(E_n\) on the input tori with
\begin{equation}
                         \|E_n\|_{a_0,b}\le b^2+\epsilon_n,
                         \qquad \epsilon_n\downarrow0.
 \label{fml-near-minimal}
\end{equation}
Average each family over the finite translations by
\(J\)-cells, translating its carriers and arguments
together. Its anchored norm does not increase, because
the translations permute the root sums. Its normalized
total remains exactly the same by the total translation
invariance of \(x_n\). The resulting array is covariant
under \(J\)-cell translations, while the total retains
all original unit-translation, point and local gauge
identities. No coefficientwise gauge invariance is needed.

Apply the proved analytic-interaction compactness argument
at the same open radius \(b\). On every fixed embedded
connected carrier, finite-dimensional Montel compactness
gives local uniform convergence of its coefficient function.
A countable diagonal extraction gives an infinite family
\(E_\infty\). Fatou on finite rooted carrier collections,
followed by radii increasing to \(b\), yields
\begin{equation}
                         \|E_\infty\|_{a_0,b}\le b^2.
 \label{fml-limit-norm}
\end{equation}
There is no analytic-radius loss.

Vacuum-normalized totals converge on compactly supported
fields. The anchored cardinality reserve makes the
large-carrier remainder uniformly small, and every fixed
carrier eventually embeds without wrapping. Local gauge
and spatial identities pass on their common compact-field
domains by the same total-continuity argument. Thus
\(E_\infty\) represents an invariant infinite class
\(x_\infty\).

Fixed-order total kernels converge locally, in particular
\begin{equation}
                K_{2,n}^{\rm in}\longrightarrow K_2(x_\infty).
 \label{fml-input-kernel-limit}
\end{equation}
They have the common exponential row and moment bounds
following from the activity norm. No numerical value of
\(M_n\) enters this input compactness argument.

\subsection{Joint extraction when the classical depth varies}
\label{fml-varying-depth-section}

The classical maps and references must be extracted jointly
with the input interactions. For every fixed depth shift
\(j\ge0\) and endpoint power \(L=R^q\), use the actual finite
data at depth \(M_nR^j\), on compatible periods obtained
from \(RN_n\) by finitely many divisions by \(R\) and \(L\).
Condition \eqref{fml-divisibility} ensures that all these
data exist for sufficiently large \(n\), and their periods
still tend to infinity.

For each fixed endpoint \(L\), an observed current output
fiber contains only \(L^4\) times a fixed number of link
and color coordinates. This dimension is independent of
the preceding refinement \(M_n\). The previously proved
classical inputs are uniform in that refinement:
\begin{itemize}
 \item a positive common holomorphic radius;
 \item strict support and cardinality bounds;
 \item complete \(C^{1,\alpha}\) bounds and the exact
       \(1/L\) pre-gauge amplitude;
 \item actual finite reference transport and minimum
       comparison identities;
 \item the controlled finite near-identity gauge comparisons.
\end{itemize}
Apply analytic compactness to the reference families and
complete-output-family compactness to each fixed-endpoint
map. One diagonal subsequence contains every depth shift
and every endpoint power in this countable collection.

Denote the resulting references by \(S_j^*\) and the maps
by \(\mathcal T_{j,q}^*\), with endpoint \(R^q\). The index
\(j\) records a relative depth shift. No assertion
\(S_j^*=S_{M_*R^j}\) for a finite \(M_*\) is made.
The uniform flat Maxwell normalization and exponential
kernel bounds give
\begin{equation}
                           \beta(S_j^*)=1.
 \label{fml-reference-normalization}
\end{equation}

The fixed-radius reference bounds, exact transport and
scalar semigroup pass to this joint system. For compact
fields the anchored tails justify passage of normalized
reference values, without passage of an extensive vacuum
energy. The scalar semigroup on every invariant input
class uses the actual finite small connecting gauges
and their endpoint estimates. These estimates are uniform
in \(M_n\); their compact extraction and the summable-gauge
argument therefore apply to this varying-depth sequence
as well. They never require the numerical label \(M_n\)
to be fixed.

Equivalently, for each fixed endpoint pair, pass the finite
comparison germ together with its actual controlled gauge
data, and then use holomorphic continuation on the
common field domain. The comparison germ can be smaller
for a large endpoint; the uniform pre-gauge domain remains
the same. The support estimate at fixed classical reserve
is uniform in all endpoint powers, and no high-jet
field-dilation radius is used here.

The resulting system satisfies
\begin{equation}
 \mathcal T_{j,q}^*S_j^*=S_{j+q}^*,\qquad
 \mathcal T_{j+q,r}^*\mathcal T_{j,q}^*
                         =\mathcal T_{j,q+r}^*.
 \label{fml-limit-semigroup}
\end{equation}
It also has the uniform quadratic-only spatial estimate
at fixed \(b\) for arbitrarily large endpoints \(R^q\).
It therefore satisfies precisely the hypotheses of
the exact infinite Maxwell-preservation proof, which gives
\begin{equation}
                   \beta(\mathcal T_{0,1}^*x)=\beta(x)
                   \quad\hbox{for every invariant }x.
 \label{fml-limit-preservation}
\end{equation}

This is a varying-depth structural limit, not an invocation
of a fixed-\(M\) limit theorem. Different limiting systems
need not coincide. Every possible system extracted this
way has \eqref{fml-limit-preservation}, which is the only
property needed below.

\subsection{The transported Hessian and its moment tails}
\label{fml-Hessian-section}

Let
\begin{equation}
                       H_n=Da_{M_n,R,N_n}(0)
 \label{fml-linear-response}
\end{equation}
be the actual linearized pre-gauge angle map. It takes
output current-angle directions on period \(N_n\) to
input current-angle directions on period \(RN_n\).
Its transpose in the scalar Hessian formula is the
ordinary counting-coordinate transpose: all physical
rescaling is already in this angle map.

The complete support-family estimate and Cauchy at zero
give
\begin{equation}
 |H_n(b,c)|
       \le C_R e^{-\kappa_R d(\operatorname{root}(b),c)}.
 \label{fml-response-envelope}
\end{equation}
The root is the final-unit cell containing the input
bond after division of its coordinates by \(R\).
Fixed internal link and color conventions are absorbed
into \(C_R\). One may retain the amplitude \(1/R\), but
\(R\) is fixed in this proof.

Each root contains only \(O(R^4)\) input coordinates.
Thus \eqref{fml-response-envelope} gives both row and
column exponential Schur bounds, with constants independent
of \(M_n\) and volume. Joint local holomorphic convergence
implies local convergence \(H_n\to H_*\).

The total first derivative of each input vanishes by
global semisimple color invariance, and the actual
pre-gauge map vanishes at zero. Consequently the exact
finite Hessian chain rule is
\begin{equation}
                       K_{2,n}^{\rm out}
                           =H_n^T K_{2,n}^{\rm in}H_n.
 \label{fml-Hessian-product}
\end{equation}
There is no term involving the second derivative of
the background map.

The exponential envelopes for the two response legs
and the input kernel give a uniform output exponential
row bound. To see this, connect the two output source
roots through the intermediate input roots. A fixed
fraction of the minimum response and input decay
exponent pays the output separation; the remaining
geometric sums are bounded by their infinite-lattice
sums. The \(O(R^4)\) fiber multiplicity is fixed. No
factor from \(M_n^4\) fine coordinates appears.

These bounds also prove local convergence of
\eqref{fml-Hessian-product}: truncate its two intermediate
root sums to a fixed neighborhood, pass the finite
sum entrywise, and then remove the truncation using
the common exponential majorant. Hence
\begin{equation}
 K_{2,*}^{\rm out}
       =H_*^T K_2(x_\infty)H_*,
 \label{fml-output-kernel-limit}
\end{equation}
which is the Hessian of
\(\mathcal T_{0,1}^*x_\infty\). The same estimates give
uniformly summable second displacement moments.

\subsection{Convergence of the actual finite Maxwell ratios}
\label{fml-moment-section}

For locally converging torus kernels with a uniform
exponential row bound, the symmetric short-lift second
moments converge to the true infinite second moment.
First restrict to a fixed displacement radius; that
finite sum converges locally. The remaining second
moment is bounded uniformly by an exponentially decaying
tail in that radius. The short-lift cutoff
\(\lfloor N_n/16\rfloor\) tends to infinity, so it
produces no residual term. Fixed midpoint shifts are
handled by the same majorant.

The finite reference denominators tend uniformly to one
by the established flat \(q_M\) periodization estimate,
including for varying \(M_n\). Therefore
\begin{align}
 \beta_{M_n,RN_n}(x_n)
                 &\longrightarrow\beta(x_\infty),\notag\\
 \beta_{M_nR,N_n}(\mathcal T_{M_n,R,N_n}x_n)
                 &\longrightarrow
                     \beta(\mathcal T_{0,1}^*x_\infty).
 \label{fml-ratio-limits}
\end{align}
These are the actual finite reference-normalized ratios.
Their full nonlinear references have not been replaced
by canonical periodizations, and no extensive total
action limit is used.

\subsection{Uniform contradiction and consequences}
\label{fml-proof-section}

\begin{proof}[Proof of Theorem~\ref{fml-theorem}]
If \eqref{fml-vanishing} failed, there would be
\(\delta>0\) and sequences as above with
\[
                     |\Lambda_{M_n,R,N_n}x_n|\ge\delta
\]
for every \(n\). Extract the joint subsequence constructed
in the preceding sections. By \eqref{fml-ratio-limits},
the limiting difference is
\[
 \beta(\mathcal T_{0,1}^*x_\infty)-\beta(x_\infty)=0
\]
by \eqref{fml-limit-preservation}. This contradicts the
lower bound and proves the theorem.
\end{proof}

On extracted inputs the conclusion reads
\begin{equation}
 \sup_{\substack{M,\ N\in\mathcal P,\ N\ge K\\
           \beta_{M,RN}(x)=0,\ \|x\|\le b^2}}
       |\beta_{MR,N}(\mathcal T_{M,R,N}x)|
                            \longrightarrow0.
 \label{fml-extracted-conclusion}
\end{equation}
The unextracted statement \eqref{fml-modulus} is slightly
stronger and needs no projection-loss factor. By linear
homogeneity it gives, for arbitrary invariant inputs,
\begin{equation}
 |\Lambda_{M,R,N}x|
       \le \omega_{R,J,b}(K)b^{-2}\|x\|,
       \qquad N\in\mathcal P,\ N\ge K.
 \label{fml-scaled-bound}
\end{equation}

For a prescribed per-step drift tolerance, choose \(K\)
so that \(\omega_{R,J,b}(K)\) is sufficiently small.
Every step of a permitted finite history whose output
period remains at least \(K\) then has the same leakage
bound, uniformly in the preceding depth and the length
of that history. On the established envelope
\(\|x\|\le C_{\rm int}b^2\), multiply the tolerance by
\(C_{\rm int}\), or choose \(K\) so that
\(\omega_{R,J,b}(K)\) is the desired tolerance divided
by this fixed constant.

\subsection{Qualifications}
\label{fml-scope-section}

The result is qualitative. It supplies no exponential
rate, no summability of leakage over infinitely many
changing periods, and no exact zero at finite volume.
The per-step consequence is a magnitude bound, not a
sign or an accumulated-drift assertion.

The cofinal divisibility assumption
\eqref{fml-divisibility} is part of the theorem. An
extension to arbitrary increasing compatible multiples
with bounded endpoint divisibility would require an
additional uniform period-extension or local-volume
comparison. It does not follow merely from \(N\to\infty\).

The limiting systems have not been identified with
one another. Exact infinite Maxwell preservation holds
in each separately and rules out every nonzero leakage
subsequence. No quantum drift sign, determinant value,
infinite physical running trajectory, or global
measure interpretation follows from this compactness
argument.
